\section{Diario della riunione}

\subsection{Dubbi sulla corretta gestione dei test di unità/integrazione del sistema LLM}
Abbiamo chiesto se ci fosse la necessità di testare le risposte del modello di LLM, anche dal punto di vista delle azioni intraprese sul database. Questo dubbio sorge in quanto
ogni risposta di un LLM consuma token della chiave fornita dall'azienda.\\
La risposta è stata che consumare i token per svolgere i test non è un problema per l'azienda, che si è anzi impegnata a fornire un set di test basato su domande frequenti
dei loro utenti.

\subsection{Dubbi requisiti desiderabili legati al caso d'uso UC\_23}
Abbiamo chiesto conferma che i requisiti legati al caso d'uso UC\_23 (fornire feedback all'AI) fossero effettivamente desiderabili 
(per maggiori informazioni consultare l'\href{https://nullpointersgroup.github.io/Documentazione/output/PB/Documenti%20Esterni/Analisi_dei_Requisiti.pdf}{Analisi dei Requisiti}). \\
La risposta è stata la conferma che questi requisiti possono essere considerati a tutti gli effetti desiderabili.

\subsection{Presentazione dello stato attuale del prodotto e richiesta di presentazione MVP in presenza}
Abbiamo presentato il prodotto allo stato attuale, concentrandoci sulle nuove funzionalità di frontend e di autenticazione.
Abbiamo inoltre chiesto la possibilità di presentare l'MVP in presenza presso la loro sede, entro due settimane dalla data odierna.\\
La risposta della azienda proponente è stata positiva.